% CV freelance - Stephane KPOVIESSI
% Working draft inspired by ModernFancyCV, intentionally simplified.
% Compile target: pdfLaTeX

\documentclass[10pt,a4paper]{article}

\usepackage[a4paper,left=0.95cm,right=0.95cm,top=0.55cm,bottom=0.55cm]{geometry}
\usepackage[utf8]{inputenc}
\usepackage[T1]{fontenc}
\usepackage{lmodern}
\usepackage{xcolor}
\usepackage{enumitem}
\usepackage[hidelinks]{hyperref}
\usepackage{tabularx}
\usepackage{titlesec}
\usepackage{fontawesome5}

\definecolor{accent}{HTML}{1F4E79}
\definecolor{accentdark}{HTML}{163A5A}
\definecolor{text}{HTML}{2F3437}
\definecolor{muted}{HTML}{666D70}
\definecolor{rule}{HTML}{D5E2EC}
\definecolor{sidebg}{HTML}{F2F7FB}

\pagestyle{empty}
\setlength{\parindent}{0pt}
\setlength{\parskip}{3.8pt}
\setlist[itemize]{leftmargin=*, topsep=3pt, itemsep=2.8pt, parsep=0.4pt, partopsep=0pt}

\titleformat{\section}
  {\fontsize{10.8pt}{12.4pt}\selectfont\bfseries\color{accent}}
  {}{0pt}{}
  [\vspace{2pt}{\color{rule}\titlerule}\vspace{3pt}]
\titlespacing*{\section}{0pt}{9pt}{4.5pt}

\newcommand{\smallsection}[2]{%
  \vspace{9pt}
  {\fontsize{8.9pt}{10.45pt}\selectfont\bfseries\color{accentdark}%
  {\scriptsize\color{accent}#1}\hspace{3.5pt}#2}\par
  \vspace{3.2pt}
}

\newcommand{\contactline}[2]{%
  \noindent\makebox[0.34cm][l]{{\color{accent}\scriptsize#1}}%
  \begin{minipage}[t]{\dimexpr\linewidth-0.36cm\relax}%
  \fontsize{8.15pt}{9.5pt}\selectfont #2%
  \end{minipage}\par
  \vspace{1pt}
}

\newcommand{\sideitem}[2]{%
  \textbf{#1}\par
  #2\par
  \vspace{5.7pt}
}

\newcommand{\certitem}[2]{%
  \textbf{#1}\par
  {\fontsize{7.15pt}{8.1pt}\selectfont\color{muted}#2}\par
  \vspace{4.4pt}
}

\newcommand{\role}[3]{%
  \vspace{2.2pt}
  \begin{tabularx}{\linewidth}{@{}Xr@{}}
  \textbf{#1} & {\color{muted}#2}
  \end{tabularx}\vspace{0.8pt}\\
  {\color{muted}#3}\par
  \vspace{2pt}
}

\newcommand{\projectline}[2]{%
  \textbf{#1} - #2\par
}
\newcommand{\githublink}[2]{%
  \par{\footnotesize\color{accent}\faIcon{github}\ \href{#1}{#2}}\par
}

\newsavebox{\maincolumnbox}
\newlength{\maincolumnheight}

\begin{document}
\fontsize{9.55pt}{11.6pt}\selectfont
\sloppy

\begin{minipage}[t]{0.52\textwidth}
{\LARGE\bfseries\color{accentdark} Stéphane KPOVIESSI}\par
\vspace{1pt}
{\large\color{accent}AI Engineer Freelance}\par
\vspace{1pt}
{\color{muted}Micro-entrepreneur \textbullet{} Disponible pour missions}
\end{minipage}
\hfill
\begin{minipage}[t]{0.42\textwidth}
\raggedleft
\vspace{2pt}
{\normalsize\bfseries\color{accent}Agents IA \textbullet{} RAG \textbullet{} Automatisation}\par
\vspace{3pt}
{\small\color{muted}
Conception d'agents IA, systèmes RAG et automatisations intégrables aux environnements métier.
}
\end{minipage}

\vspace{6pt}
{\color{rule}\hrule height 1pt}
\vspace{7pt}

\begin{lrbox}{\maincolumnbox}
\begin{minipage}[t]{0.695\textwidth}
\vspace{0pt}
\raggedright

\section{AI Engineering \& Automatisation}
J'aide les entreprises à identifier leurs besoins d'automatisation, d'outillage interne et d'intégration IA, puis à les concrétiser en workflows, outils open source, assistants, systèmes RAG et agents IA. Du cadrage au POC, puis à l'intégration métier.

\textbf{\color{accent}Processus métier -> Opportunité -> POC / Solution -> Intégration}

\vspace{7pt}
\section{Seshat}
\role{Fondateur \& AI Engineer}{Sept. 2026 - Aujourd'hui}{Plateforme open source d'IA agentique et d'automatisation}
{\footnotesize \href{https://seshat-ai.com}{seshat-ai.com} \textbullet{} \href{https://github.com/KPO-Tech/seshat}{github.com/KPO-Tech/seshat}}\par
\begin{itemize}
  \item Conception d'une plateforme d'IA agentique et d'automatisation pour créer des agents et workflows intégrables aux environnements d'entreprise.
  \item Architecture Python/Go, services backend local/cloud et application desktop Electron/React.
  \item Orchestration d'agents : tool calling, MCP, human-in-the-loop, tâches planifiées et workflows.
  \item RAG hybride et multi-LLM : ingestion documentaire, embeddings, PostgreSQL/pgvector, Chroma, OpenAI, Mistral, Gemini et Ollama.
\end{itemize}
{\footnotesize\textbf{Entreprise :} permissions, auditabilité, sécurité. \quad \textbf{Usages :} assistants, RAG, automatisation.}

\section{Expériences}
\role{Allianz France - Ingénieur Data / Data Scientist}{Nov. 2025 - Avr. 2026}{Paris, France}
\begin{itemize}
  \item Migration d'un datamart Assurance de SAS vers Python/PySpark, avec validation de parité fonctionnelle sur plus de 15 000 lignes de code.
  \item Architecture Médaillon sur Azure Databricks intégrant plus de 150 règles de gestion métier.
  \item Pipeline de Data Quality géospatiale et framework de tests automatisés SAS/Python pour validation de KPI et non-régression.
\end{itemize}

\vspace{4pt}
\role{LYNE ET FRERES SARL - Data Analyst}{2023 - 2025}{Missions d'appui stratégique - Bénin / À distance}
\begin{itemize}
  \item Analyse de données commerciales, opérationnelles et agricoles avec Python/Pandas pour suivre les KPI et soutenir les décisions.
  \item Création de rapports, tableaux de bord et interfaces Streamlit / web pour restituer les analyses aux utilisateurs métier.
  \item Déploiement d'automatisations et d'outils IA internes open source (OpenWebUI, modèles ouverts), avec accompagnement des employés.
\end{itemize}

\vspace{4pt}
\role{Groupe Sylvagreg - Machine Learning Engineer}{Mai 2024 - Août 2024}{Lille, France}
R\&D computer vision pour la reconstruction 3D à partir d'images 2D : modèles PyTorch, pipelines de prétraitement/inférence, Docker, AWS et visualisation métier.\par
{\fontsize{7.5pt}{9pt}\selectfont\color{accent}\faIcon{github}\ \href{https://github.com/EngineerProjects/Sylva3dGUI}{github.com/EngineerProjects/Sylva3dGUI}\enspace\textbullet\enspace\href{https://github.com/EngineerProjects/sylva3D}{github.com/EngineerProjects/sylva3D}}\par

\section{Projets techniques}
\projectline{TECH WATCH AGENT}{LangGraph, RAG, FastAPI, pgvector, Docker}
Plateforme de veille multi-agents : planification de sujets, collecte multi-source, extraction d'éléments clés, synthèse sourcée, rapports ciblés, workflows LangGraph et mémoire vectorielle pgvector.
\githublink{https://github.com/EngineerProjects/tech-watch-agent}{github.com/EngineerProjects/tech-watch-agent}

\vspace{4pt}
\projectline{BI RETAIL SOLUTION}{PySpark, Airflow, PostgreSQL, Power BI, Docker}
Pipeline Data ventes/supply chain : ingestion, transformations PySpark, orchestration Airflow, DWH PostgreSQL et tableaux de bord Power BI.
\githublink{https://github.com/EngineerProjects/BI_Retail}{github.com/EngineerProjects/BI\_Retail}

\end{minipage}
\end{lrbox}
\setlength{\maincolumnheight}{\dimexpr\ht\maincolumnbox+\dp\maincolumnbox\relax}

\begin{minipage}[t]{0.285\textwidth}
\vspace{0pt}

\hspace*{-0.18cm}%
\colorbox{sidebg}{%
  \hspace*{0.18cm}%
  \begin{minipage}[t][\maincolumnheight][t]{0.91\linewidth}
  \vspace{8pt}
  \fontsize{8.45pt}{10.35pt}\selectfont
  \raggedright

  \smallsection{\faIcon{address-card}}{Contact}
  \contactline{\faIcon{map-marker-alt}}{Cergy / Île-de-France}
  \contactline{\faIcon{phone}}{+33 7 48 59 47 13}
  \contactline{\faIcon{envelope}}{\href{mailto:oastephanekpoviessi@gmail.com}{oastephanekpoviessi@gmail.com}}
  \contactline{\faIcon{link}}{\href{https://linkedin.com/in/stephanekpoviessi}{linkedin.com/in/stephanekpoviessi}}
  \contactline{\faIcon{globe}}{\href{https://seshat-ai.com}{seshat-ai.com}}
  \contactline{\faIcon{code-branch}}{\href{https://github.com/KPO-Tech}{github.com/KPO-Tech}}

  \smallsection{\faIcon{compass}}{Interventions}
  \sideitem{Cadrage \& cas d'usage IA}{Processus, opportunités IA, priorisation, cadrage de pilote.}
  \sideitem{Agents IA}{Orchestration, tool calling, MCP, human-in-the-loop.}
  \sideitem{RAG \& Knowledge}{Ingestion, embeddings, bases vectorielles, recherche hybride.}
  \sideitem{Automatisation}{APIs, outils internes, tâches planifiées, workflows métier.}

  \smallsection{\faIcon{layer-group}}{Stack}
  \sideitem{AI \& Agents}{LLM, RAG, LangGraph, MCP, Tool Calling.}
  \sideitem{Backend \& Data}{Python, Go, FastAPI, PostgreSQL, pgvector, PySpark.}
  \sideitem{Cloud \& Engineering}{Azure, Databricks, Airflow, Docker, Git/CI-CD.}
  \sideitem{Frontend}{React, Electron.}

  \smallsection{\faIcon{graduation-cap}}{Formation}
  \sideitem{JUNIA ISEN Lille}{Diplôme d'Ingénieur - Big Data \& Data Science, 2026.}
  \sideitem{JUNIA ISEN Rabat}{Licence - Électronique \& Systèmes Numériques, 2023.}

  \smallsection{\faIcon{certificate}}{Certifications}
  \certitem{\href{https://learn.microsoft.com/fr-fr/training/career-paths/ai-engineer}{Microsoft AI Engineer Career Path}}{Microsoft Learn \textbullet{} \textbf{En cours}}
  \certitem{\href{https://www.coursera.org/account/accomplishments/verify/LU2ZQ972BWNF}{Create ML Models in Microsoft Azure}}{Microsoft / Coursera \textbullet{} 2023}
  \certitem{\href{https://www.coursera.org/account/accomplishments/verify/MWJS4Q3S2C3L}{Supervised ML: Regression \& Classification}}{DeepLearning.AI / Stanford Online\\2023}

  \smallsection{\faIcon{language}}{Langues}
  \textbf{Français} - langue maternelle\par
  \textbf{Anglais} - B2, Linguaskill Business
  \vspace{6pt}
  \end{minipage}
}

\end{minipage}
\hfill
\usebox{\maincolumnbox}

\end{document}
